\documentclass[book.tex]{subfiles}
\begin{document}

\chapter{Graph Networks in High Energy Physics}

\label{chap:gnn_hep}
\noindent\rule{11cm}{0.4pt} \\
\textbf{Prerequisites:} \autoref{chap:graphconvs}, \autoref{chap:equivariant_networks}, \autoref{chap:molecular_dynamics},  \autoref{chap:dft} \\
\textbf{Difficulty Level:} ** \\
\noindent\rule{11cm}{0.4pt}
\vspace{5mm}

Graph neural networks can be used for the tasks of reconstruction, tagging, generation, and end-to-end analysis.

In particular, graph neural networks have found a niche in analyzing particle accelerator data to perform useful tasks like identifying jets. Architectures for this task need to consider useful physical properties of jets which may be used to narrow down those phenomena which merit deeper consideration.

\section{Reconstruction and Identification}

Raw data from particle accelerators can take the form of a point cloud. A graph of relationships between these events has to be reconstructed. Traditionally this is done with heuristic methods, but this can be done with graph networks.

Reconstruction approaches often build on a clustering task and may follow the form of edge strength prediction which predicts the weight of a node to node edge. Calorimeter applications focus on segmentation, with tasks like partial assignment of hits to showers.

Identification tasks focus on graph classification or segmentation. A common task is jet classification, which seeks to identify the particle that triggered a jet. Graph networks have proven useful in identifying jets and reconstructing their dynamics from the raw data.

\section{Generative Modeling}

Generative models are useful for creating simulated data since running lattice cluster simulations can be very expensive.


\end{document}